% Generated draft: requires mathematical and source review.
\documentclass[11pt]{article}
\usepackage[margin=25mm]{geometry}
\usepackage{fontspec}
\setmainfont{Latin Modern Roman}
\usepackage{amsmath,amssymb,mathtools,unicode-math}
\setmathfont{Latin Modern Math}
\usepackage{enumitem,needspace,xurl,hyperref}
\setlength{\parindent}{0pt}
\setlength{\parskip}{5pt}
\setlength{\emergencystretch}{3em}
\newcommand{\N}{\mathbb{N}}
\newcommand{\Z}{\mathbb{Z}}
\newcommand{\Q}{\mathbb{Q}}
\newcommand{\R}{\mathbb{R}}
\newcommand{\C}{\mathbb{C}}
\newcommand{\F}{\mathbb{F}}
\newcommand{\A}{\mathbb{A}}
\newcommand{\PP}{\mathbb{P}}
\newcommand{\E}{\mathbb{E}}
\newcommand{\eps}{\varepsilon}
\newcommand{\dd}{\,\mathrm{d}}
\begin{document}
\section*{UnsolvedMath ID: 30005741}
\textbf{Source problem code:} catalogue entry
\subsection{Definitions and mathematical statement}
Let \(K\) be a field of characteristic \(p>0\), and let \(L\) be a Lie algebra over \(K\). No finite-dimensionality, perfectness of \(K\), or restricted Lie algebra structure is assumed in the supplied statement. Interpret the notation \(s(L)\) as the truncated symmetric Poisson algebra
\[
A=s(L):=S_K(L)/\langle x^p:x\in L\rangle.
\]
Here \(S_K(L)\) is the symmetric algebra, equipped with the Poisson bracket extending the Lie bracket on \(L\) by bilinearity and the Leibniz rule. The displayed ideal is a Poisson ideal, so the bracket descends to \(A\). This interpretation of \(s(L)\) and the series conventions below are standard reconstructions; they are not explicitly defined in the supplied source excerpt and require comparison with the cited report.

For \(K\)-linear subspaces \(U,V\subseteq A\), let \(\{U,V\}\) denote the \(K\)-linear span of all brackets \(\{u,v\}\), and let \(UA\) denote the span of all products \(ua\). Define the ordinary and strong lower central series by
\[
\gamma_1(A)=A,\qquad
\gamma_{n+1}(A)=\{\gamma_n(A),A\},
\]
\[
\Gamma_1(A)=A,\qquad
\Gamma_{n+1}(A)=\{\Gamma_n(A),A\}A
\quad(n\geq1).
\]
Define their nilpotency classes by
\[
c(A)=\min\{c\geq1:\gamma_{c+1}(A)=0\},\qquad
c_{\mathrm{str}}(A)=\min\{c\geq1:\Gamma_{c+1}(A)=0\}.
\]
An empty minimum is assigned the value \(\infty\). Thus \(A\) is Lie nilpotent when \(c(A)<\infty\), and strongly Lie nilpotent when \(c_{\mathrm{str}}(A)<\infty\).

Define the ordinary and strong derived series by
\[
D_0(A)=A,\qquad D_{n+1}(A)=\{D_n(A),D_n(A)\},
\]
\[
E_0(A)=A,\qquad E_{n+1}(A)=\{E_n(A),E_n(A)\}A
\quad(n\geq0).
\]
Their derived lengths are
\[
\ell(A)=\min\{d\geq0:D_d(A)=0\},\qquad
\ell_{\mathrm{str}}(A)=\min\{d\geq0:E_d(A)=0\},
\]
again assigning \(\infty\) to an empty minimum. Solvability means \(\ell(A)<\infty\), and strong solvability means \(\ell_{\mathrm{str}}(A)<\infty\).

\begin{enumerate}
\item For every field \(K\) of characteristic \(p\in\{2,3\}\) and every Lie algebra \(L\) over \(K\), does
\[
c(s(L))<\infty\quad\Longrightarrow\quad
c_{\mathrm{str}}(s(L))=c(s(L))
\]
hold?
\item For every field \(K\) of characteristic \(p>2\) and every Lie algebra \(L\) over \(K\), does
\[
\ell(s(L))<\infty\quad\Longrightarrow\quad
\ell_{\mathrm{str}}(s(L))=\ell(s(L))
\]
hold?
\end{enumerate}
The second question is included because it appears in the supplied reconstructed statement. The supplied verbatim excerpt only asks whether an unspecified preceding assertion remains true in characteristics \(2,3\); it does not independently establish the scope of either question.
\subsection{Short English statement}
For the truncated symmetric Poisson algebra of an arbitrary Lie algebra, ask whether Lie nilpotency forces the same strong Lie nilpotency class in characteristics \(2\) and \(3\), and whether solvability forces the same strong derived length in every characteristic greater than \(2\). The supplied excerpt lacks its preceding context, so the reconstructed scope and definitions require human review.
\subsection{Sources}
\begin{itemize}
\item Supplied local UnsolvedMath registry record, permanent ID \(30005741\), source code \texttt{catalogue entry}; collection: Oberwolfach Reports Open Problems. Registry locations identified by the project instructions are \texttt{data/unsolvedmath/problems.json} and the supplementary \texttt{catalog/problems.json}; the supplied record does not specify which contains this entry.
\item Supplied report citation: \emph{Mini-Workshop: Poisson and Poisson-type algebras (2024)}, Oberwolfach Reports, pp.\ 2681--2715, DOI \texttt{10.4171/owr/2023/46}. \url{https://doi.org/10.4171/owr/2023/46}. The literature-source label is \emph{Oberwolfach Report 46/2023}; the record gives proposed year 2023. The supplied original excerpt labels the question Problem 1 but omits the antecedent of ``this.'' No individual proposer is supplied.
\item Monteiro Alves and Petrogradsky, \emph{Lie structure of truncated symmetric Poisson algebras}. \url{https://doi.org/10.1016/j.jalgebra.2017.05.035}. This reference and attribution are supplied by the record.
\item Literature assessment dated exactly \texttt{2026-08-22}: the supplied curation labels the problem open and states that the 2017 theorem gives equality of Lie and strong Lie nilpotency classes only for \(p\geq5\), while the 2023 report leaves \(p=2,3\) as Problem 1. This is supplied dated curation, not an independently verified literature assessment. The background expressly requests independent verification of time-sensitive status and cited claims.
\item The registry labels the statement \texttt{corrected\_verified} and reports that it was reconstructed and checked against the official Oberwolfach report. These are supplied provenance claims; this draft does not independently verify them or claim to have visited either external source.
\end{itemize}
\end{document}
